\documentclass{beamer}
\usepackage[utf8]{inputenc}

\usetheme{Madrid}
\usecolortheme{default}
\usepackage{amsmath,amssymb,amsfonts,amsthm}
\usepackage{mathtools}
\usepackage{txfonts}
\usepackage{tkz-euclide}
\usepackage{listings}
\usepackage{adjustbox}
\usepackage{array}
\usepackage{gensymb}
\usepackage{tabularx}
\usepackage{gvv}
\usepackage{lmodern}
\usepackage{circuitikz}
\usepackage{tikz}
\lstset{literate={·}{{$\cdot$}}1 {λ}{{$\lambda$}}1 {→}{{$\to$}}1}
\usepackage{graphicx}

\setbeamertemplate{page number in head/foot}[totalframenumber]

\usepackage{tcolorbox}
\tcbuselibrary{minted,breakable,xparse,skins}



\definecolor{bg}{gray}{0.95}
\DeclareTCBListing{mintedbox}{O{}m!O{}}{%
  breakable=true,
  listing engine=minted,
  listing only,
  minted language=#2,
  minted style=default,
  minted options={%
    linenos,
    gobble=0,
    breaklines=true,
    breakafter=,,
    fontsize=\small,
    numbersep=8pt,
    #1},
  boxsep=0pt,
  left skip=0pt,
  right skip=0pt,
  left=25pt,
  right=0pt,
  top=3pt,
  bottom=3pt,
  arc=5pt,
  leftrule=0pt,
  rightrule=0pt,
  bottomrule=2pt,
  toprule=2pt,
  colback=bg,
  colframe=orange!70,
  enhanced,
  overlay={%
    \begin{tcbclipinterior}
    \fill[orange!20!white] (frame.south west) rectangle ([xshift=20pt]frame.north west);
    \end{tcbclipinterior}},
  #3,
}
\lstset{
    language=C,
    basicstyle=\ttfamily\small,
    keywordstyle=\color{blue},
    stringstyle=\color{orange},
    commentstyle=\color{green!60!black},
    numbers=left,
    numberstyle=\tiny\color{gray},
    breaklines=true,
    showstringspaces=false,
}

\title{5.2.19}
\date{October 4, 2025}
\author{Taraka Abhinav - EE25BTECH11016}

\begin{document}

\frame{\titlepage}
\begin{frame}{Question}
Solve the system of equations:
\begin{align}
x + y &= 14 \\
x - y &= 4
\end{align}
\end{frame}

\begin{frame}{Line Representation}
The equation of a line is
\begin{align}
\vec{n^T}\vec{x} = c
\end{align}

Line L1:
\begin{align}
\myvec{1 & 1}\myvec{x \\ y}=14
\end{align}
Line L2:
\begin{align}
\myvec{1 & -1}\myvec{x \\ y}=4
\end{align}
\end{frame}
\begin{frame}{Matrix Form}
These can be combined into the matrix form $A\vec{x}=\vec{b}$:
\begin{align}
\myvec{1 & 1 \\ 1 & -1}\myvec{x \\ y} = \myvec{14 \\ 4}
\end{align}

The following augmented matrix can be solved by Gaussian elimination 
\begin{align}
\augvec{2}{1}{1 & 1 & 14 \\ 1 & -1 & 4} 
\xleftrightarrow{R_2 \leftarrow R_2 - R_1} 
\augvec{2}{1}{1 & 1 & 14 \\ 0 & -2 & -10}
\end{align}
\end{frame}

\begin{frame}{Reduced Equations and Solution}
The matrix is in row echelon form with two non-zero rows (Rank = 2). We can use back substitution. The second row gives:
\begin{align}
-2y = -10 \implies y = 5
\end{align}
Substituting $y=5$ into the first row equation, $x+y=14$:
\begin{align}
x + 5 = 14 \implies x = 9
\end{align}
\end{frame}

\begin{frame}{Final Solution}
The unique solution is:
\begin{align}
\myvec{x \\ y} = \myvec{9 \\ 5}
\end{align}

\textbf{Conclusion:} 
The system has a \textbf{unique solution}.
\end{frame}

\begin{frame}[fragile]
    \frametitle{C Code}
    \begin{lstlisting}
#include <stdio.h>

typedef struct {
    double x, y;
} Point;

Point solve_linear_system() {
    // Solving manually for the given system
    Point p;
    // x + y = 14
    // x - y = 4
    p.x = (14 + 4)/2; // 9
    p.y = 14 - p.x;   // 5
    return p;
}

int main() {
    Point P = solve_linear_system();
    printf("Intersection Point: P(%g, %g)\n", P.x, P.y);
    return 0;
}
    \end{lstlisting}
\end{frame}

\begin{frame}[fragile]
    \frametitle{Python + C}
    \begin{lstlisting}
    import ctypes
import numpy as np
import matplotlib.pyplot as plt

class Point(ctypes.Structure):
    _fields_ = [("x", ctypes.c_double), ("y", ctypes.c_double)]

lib = ctypes.CDLL("./solve_linear.so")
lib.solve_linear_system.restype = Point

P = lib.solve_linear_system()

# Plotting
x_vals = np.linspace(0, 15, 300)
y_L = 14 - x_vals
y_K = x_vals - 4
    \end{lstlisting}
\end{frame}

\begin{frame}[fragile]
    \frametitle{Python + C}
    \begin{lstlisting}
plt.figure()
plt.plot(x_vals, y_L, label='L: x+y=14')
plt.plot(x_vals, y_K, label='K: x-y=4')

plt.scatter(P.x, P.y, color='red')
plt.text(P.x+0.3, P.y+0.3, f'P({int(P.x)},{int(P.y)})', color='red')

plt.xlim(0,15)
plt.ylim(0,15)
plt.xlabel('x')
plt.ylabel('y')
plt.title('Solution of Linear System (C + Python)')
plt.grid(True)
plt.legend()
plt.show()
    \end{lstlisting}
\end{frame}

\begin{frame}[fragile]
    \frametitle{Python Code}
    \begin{lstlisting}
import numpy as np
import matplotlib.pyplot as plt

# Coefficients for lines
# Line L: x + y = 14
# Line K: x - y = 4

# Define x range for plotting
x_vals = np.linspace(0, 15, 300)

# Line L: y = 14 - x
y_L = 14 - x_vals

# Line K: y = x - 4
y_K = x_vals - 4
    \end{lstlisting}
\end{frame}

\begin{frame}[fragile]
    \frametitle{Python Code}
    \begin{lstlisting}
# Intersection point (solving manually or using numpy)
A = np.array([[1, 1], [1, -1]])
b = np.array([14, 4])
sol = np.linalg.solve(A, b)
x_int, y_int = sol

# Plot lines
plt.figure()
plt.plot(x_vals, y_L, label='L: x+y=14')
plt.plot(x_vals, y_K, label='K: x-y=4')

# Mark intersection
plt.scatter(x_int, y_int, color='red')
plt.text(x_int+0.3, y_int+0.3, f'P({int(x_int)},{int(y_int)})', color='red')

plt.xlim(0,15)
plt.ylim(0,15)
plt.xlabel('x')
plt.ylabel('y')
plt.title('Solution of Linear System')
plt.grid(True)
plt.legend()
plt.show()
    \end{lstlisting}
\end{frame}
\begin{frame}{Plot}
\begin{figure}[h!]
    \centering
    \includegraphics[height=0.8\textheight, keepaspectratio]{figs/Figure_1.png}
    \label{figure_1}
\end{figure}
\end{frame}
\end{document}